% !TeX program = xelatex
% 更新日期：2026-11-12；含连续两个正面停止的生成函数例子。
% Overleaf 单文件版本：将本文件全部内容替换到 main.tex，编译器选择 XeLaTeX。
% 所有正文与样式均已内嵌，不需要上传其他文件。
\documentclass[UTF8,a4paper,11pt,fontset=fandol]{ctexart}
\usepackage[left=29mm,right=25mm,top=23mm,bottom=24mm]{geometry}
\usepackage{amsmath,amssymb,amsthm}
\usepackage{xcolor}
\usepackage[most]{tcolorbox}
\definecolor{pnRed}{HTML}{A62930}
\definecolor{pnBlue}{HTML}{253C68}
\definecolor{pnGold}{HTML}{BD8D24}
\definecolor{pnDefinitionBg}{HTML}{F0E3D7}
\definecolor{pnPropositionBg}{HTML}{F8F5EF}
\definecolor{pnExampleBg}{HTML}{EFE5CB}
\newcommand{\pnE}{\mathbb E}
\newcommand{\pnVar}{\operatorname{Var}}
\newcommand{\pnCov}{\operatorname{Cov}}
\newcommand{\pnIndicator}{\mathbf 1}
\newcounter{pnstatement}[section]
\renewcommand{\thepnstatement}{\thesection.\arabic{pnstatement}}
\tcbset{pnbase/.style={enhanced,breakable,sharp corners,boxrule=0pt,
 borderline west={1.2pt}{0pt}{pnBlue},left=9pt,right=9pt,top=7pt,bottom=7pt,
 before skip=10pt,after skip=10pt,fonttitle=\bfseries,coltitle=pnBlue,
 attach title to upper={\quad},title after break={},
 overlay unbroken and first={\fill[tcbcolframe] ([xshift=-9pt,yshift=-9pt]frame.north west) rectangle ++(4pt,4pt);}}}
\newenvironment{pnDefinition}[1]{\refstepcounter{pnstatement}\begin{tcolorbox}[
 pnbase,colback=pnDefinitionBg,colframe=pnRed,coltitle=pnRed,
 borderline west={1.2pt}{0pt}{pnRed},title={定义~\thepnstatement\ (#1).}]}{\end{tcolorbox}}
\newenvironment{pnProposition}[1]{\refstepcounter{pnstatement}\begin{tcolorbox}[
 pnbase,colback=pnPropositionBg,colframe=pnBlue,
 title={命题~\thepnstatement\ (#1).}]}{\end{tcolorbox}}
\newenvironment{pnExample}[1]{\refstepcounter{pnstatement}\begin{tcolorbox}[
 pnbase,colback=pnExampleBg,colframe=pnGold,coltitle=pnGold,
 borderline west={1.2pt}{0pt}{pnGold},title={例~\thepnstatement\ (#1).}]}{\end{tcolorbox}}
\newenvironment{pnRemark}[1]{\par\medskip\noindent
 \refstepcounter{pnstatement}\textcolor{pnRed}{注~\thepnstatement\ (#1).}\ }{\par\medskip}
\renewcommand{\proofname}{证明}
\usepackage{fancyhdr}
\usepackage{tikz}
\usepackage{eso-pic}
\usepackage[unicode,hidelinks]{hyperref}
\ctexset{section={format=\Large\bfseries\color{pnRed}},
 subsection={format=\large\bfseries\color{pnRed}}}
\setlength{\parskip}{0.4em}
\setlength{\headheight}{14pt}
\pagestyle{fancy}
\fancyhf{}
\renewcommand{\headrulewidth}{0pt}
\fancyfoot[L]{\small\color{gray}离散数学与结构（I）}
\fancyfoot[R]{\small\color{gray}\thepage}
\AddToShipoutPictureBG{\begin{tikzpicture}[remember picture,overlay]
 \fill[pnExampleBg!55] ([xshift=10.5mm,yshift=-19mm]current page.north west)
 rectangle ([xshift=15mm,yshift=20mm]current page.south west);
 \foreach \x in {10.5,12,15}{\draw[pnGold!55,line width=0.3pt]
 ([xshift=\x mm,yshift=-19mm]current page.north west) --
 ([xshift=\x mm,yshift=20mm]current page.south west);}
\end{tikzpicture}}
\begin{document}
% 接续现有讲义，从第 9 节开始；若希望从第 1 节开始，把下面的 8 改为 0。
\setcounter{section}{8}
\section{概率与概率生成函数}
\label{sec:pn-probability}

本节讨论离散概率。统一记 $\mathbb N=\{0,1,2,\ldots\}$，
$\mathbb N_{>0}=\{1,2,\ldots\}$；$p(\omega)$ 表示样本点的概率，
$\Pr[A]$ 表示事件的概率，$p_X(x)$ 表示随机变量的概率质量函数。

\subsection{概率空间、条件概率与独立性}

\begin{pnDefinition}{离散概率空间与事件}
设 $\Omega$ 是有限或可数的样本空间，函数 $p\colon\Omega\to\mathbb R$ 满足
\[
 p(\omega)\geq 0\quad(\omega\in\Omega),
 \qquad \sum_{\omega\in\Omega}p(\omega)=1.
\]
称 $(\Omega,p)$ 为离散概率空间。子集 $A\subseteq\Omega$ 称为事件，其概率定义为
\[
 \Pr[A]=\sum_{\omega\in A}p(\omega).
\]
\end{pnDefinition}

例如，掷一枚骰子可以取 $\Omega=\{1,2,3,4,5,6\}$。
对任意事件 $A,B$，有
\[
 \Pr[A\cup B]=\Pr[A]+\Pr[B]-\Pr[A\cap B]
 \leq \Pr[A]+\Pr[B].
\]

\begin{pnDefinition}{条件概率与条件概率空间}
设 $\Pr[B]>0$，定义
\[
 \Pr[A\mid B]=\frac{\Pr[A\cap B]}{\Pr[B]}.
\]
也可以在原样本空间上定义新的点概率
\[
 p_B(\omega)=
 \begin{cases}
  p(\omega)/\Pr[B],&\omega\in B,\\
  0,&\omega\notin B.
 \end{cases}
\]
于是 $(\Omega,p_B)$ 是概率空间，且 $\Pr[A\mid B]=\sum_{\omega\in A}p_B(\omega)$。
\end{pnDefinition}

\begin{pnProposition}{Bayes 公式}
若 $\Pr[A]>0$ 且 $\Pr[B]>0$，则
\[
 \Pr[A\mid B]=\frac{\Pr[A]\Pr[B\mid A]}{\Pr[B]}.
\]
\end{pnProposition}
\begin{proof}
由条件概率的定义，
\[
 \Pr[A\mid B]\Pr[B]=\Pr[A\cap B]
 =\Pr[B\mid A]\Pr[A].\qedhere
\]
\end{proof}

\begin{pnExample}{疾病检测}
设 $A$ 为患病事件，$B$ 为检测呈阳性事件，且
\[
 \Pr[B\mid A]=99\%,\qquad \Pr[\overline B\mid\overline A]=99\%.
\]
若患病率为 $\Pr[A]=0.1\%$，则
\begin{align*}
 \Pr[B]
 &=\Pr[B\mid A]\Pr[A]
   +\Pr[B\mid\overline A]\Pr[\overline A]\\
 &=0.99\times0.001+0.01\times0.999=0.01098,\\
 \Pr[A\mid B]
 &=\frac{0.99\times0.001}{0.01098}
 =\frac{11}{122}\approx 9.02\%\approx\frac1{11}.
\end{align*}
\end{pnExample}

\begin{pnDefinition}{事件的独立性}
若 $\Pr[A\cap B]=\Pr[A]\Pr[B]$，称事件 $A,B$ 独立。
当 $\Pr[B]>0$ 时，等价地有 $\Pr[A\mid B]=\Pr[A]$。
若进一步 $\Pr[\overline B]>0$，则也有 $\Pr[A\mid\overline B]=\Pr[A]$。
\end{pnDefinition}

\subsection{随机变量与分布}

\begin{pnDefinition}{随机变量与分布}
随机变量是定义在样本空间上的函数 $X\colon\Omega\to S$。
它在 $S$ 上诱导的分布为
\[
 p_X(x)=\Pr[X=x]
 =\sum_{\substack{\omega\in\Omega\\X(\omega)=x}}p(\omega).
\]
因此也可以把 $(S,p_X)$ 看作概率空间。
\end{pnDefinition}

\begin{pnExample}{抛硬币与指示随机变量}
抛 $n$ 次硬币，令 $\Omega=\{H,T\}^n$，定义
\[
 X_i(\omega)=
 \begin{cases}1,&\omega_i=H,\\0,&\omega_i=T,\end{cases}
 \qquad 1\leq i\leq n.
\]
则 $X=\sum_{i=1}^nX_i$ 为正面出现的总次数，
$Y=\min\{i:X_i=1\}$ 为第一次出现正面的位置。
为使后一个定义在整个 $\Omega$ 上有意义，约定未出现正面时 $Y=+\infty$。

一般地，事件 $A$ 的指示随机变量定义为
\[
 \pnIndicator_A\colon\Omega\to\{0,1\},\qquad
 \pnIndicator_A(\omega)=
 \begin{cases}1,&\omega\in A,\\0,&\omega\notin A.\end{cases}
\]
\end{pnExample}

\begin{pnDefinition}{联合分布、边缘分布与条件分布}
设 $X,Y$ 为随机变量，联合分布定义为
\[
 p_{X,Y}(x,y)=\Pr[X=x,Y=y].
\]
其边缘分布满足
\[
 p_X(x)=\sum_y p_{X,Y}(x,y).
\]
当 $p_Y(y)>0$ 时，条件分布定义为
\[
 p_{X\mid Y}(x\mid y)=\Pr[X=x\mid Y=y]
 =\frac{p_{X,Y}(x,y)}{p_Y(y)}.
\]
\end{pnDefinition}

因此，在 $p_Y(y)>0$ 时有
\[
 p_{X,Y}(x,y)=p_{X\mid Y}(x\mid y)p_Y(y),
 \qquad
 p_X(x)=\sum_{y:\,p_Y(y)>0}p_{X\mid Y}(x\mid y)p_Y(y).
\]
这里的求和带有权重 $p_Y(y)$，不能直接把条件分布对 $y$ 相加。
若 $p_Y(y)=0$，则所有 $p_{X,Y}(x,y)=0$，该项不影响边缘分布。

\begin{pnDefinition}{随机变量的独立性}
若对所有 $x,y$ 都有
\[
 p_{X,Y}(x,y)=p_X(x)p_Y(y),
\]
则称 $X,Y$ 独立。
\end{pnDefinition}

\subsection{期望与条件期望}

本小节的期望针对实值随机变量；涉及带符号的无穷求和时，假设所需的绝对期望有限。

\begin{pnDefinition}{期望}
随机变量 $X$ 的期望定义为
\[
 \pnE[X]=\sum_x p_X(x)x
 =\sum_{\omega\in\Omega}p(\omega)X(\omega).
\]
\end{pnDefinition}

\begin{pnProposition}{期望的线性性}
对可积的随机变量 $X,Y$，有 $\pnE[X+Y]=\pnE[X]+\pnE[Y]$。
\end{pnProposition}
\begin{proof}
由样本空间上的表达式，
\begin{align*}
 \pnE[X+Y]
 &=\sum_\omega p(\omega)\bigl(X(\omega)+Y(\omega)\bigr)\\
 &=\sum_\omega p(\omega)X(\omega)
   +\sum_\omega p(\omega)Y(\omega)
 =\pnE[X]+\pnE[Y].
\end{align*}
\end{proof}

\begin{pnDefinition}{关于事件的条件期望}
若 $\Pr[A]>0$，定义
\[
 \pnE[X\mid A]
 =\sum_x\Pr[X=x\mid A]x
 =\sum_\omega p_A(\omega)X(\omega).
\]
\end{pnDefinition}

\begin{pnDefinition}{关于随机变量的条件期望}
对 $p_Y(y)>0$ 的 $y$，令
\[
 f(y)=\pnE[X\mid Y=y]
 =\sum_x p_{X\mid Y}(x\mid y)x.
\]
定义随机变量 $\pnE[X\mid Y]=f(Y)$，即
\[
 \bigl(\pnE[X\mid Y]\bigr)(\omega)=f\bigl(Y(\omega)\bigr).
\]
在 $p_Y(y)=0$ 处可以任意指定 $f(y)$，例如取 $0$。
\end{pnDefinition}

固定 $y$ 后，$\pnE[X\mid Y=y]$ 是一个数；$\pnE[X\mid Y]$ 则是随机变量。

\begin{pnProposition}{重期望公式}
若 $X$ 可积，则 $\pnE[X]=\pnE\bigl[\pnE[X\mid Y]\bigr]$。
\end{pnProposition}
\begin{proof}
以下只对 $p_Y(y)>0$ 的 $y$ 求和：
\begin{align*}
 \pnE\bigl[\pnE[X\mid Y]\bigr]
 &=\sum_y p_Y(y)\pnE[X\mid Y=y]\\
 &=\sum_y\sum_x p_Y(y)p_{X\mid Y}(x\mid y)x\\
 &=\sum_x\sum_y p_{X,Y}(x,y)x
 =\sum_x p_X(x)x=\pnE[X].
\end{align*}
\end{proof}

\begin{pnProposition}{独立随机变量乘积的期望}
若 $X,Y$ 独立且可积，则 $\pnE[XY]=\pnE[X]\pnE[Y]$。
\end{pnProposition}
\begin{proof}
\begin{align*}
 \pnE[XY]
 &=\sum_{x,y}p_{X,Y}(x,y)xy
 =\sum_{x,y}p_X(x)p_Y(y)xy\\
 &=\left(\sum_xp_X(x)x\right)
   \left(\sum_yp_Y(y)y\right)
 =\pnE[X]\pnE[Y].
\end{align*}
\end{proof}

\subsection{方差与协方差}

本小节假设涉及的随机变量具有有限二阶矩。

\begin{pnDefinition}{方差、协方差与协方差矩阵}
随机变量 $X$ 的方差定义为
\[
 \pnVar(X)=\pnE\bigl[(X-\pnE[X])^2\bigr].
\]
随机变量 $X,Y$ 的协方差定义为
\[
 \pnCov(X,Y)=\pnE\bigl[(X-\pnE[X])(Y-\pnE[Y])\bigr],
 \qquad \pnVar(X)=\pnCov(X,X).
\]
二元随机向量的协方差矩阵为
\begin{align*}
 &\begin{pmatrix}
  \pnVar(X)&\pnCov(X,Y)\\
  \pnCov(Y,X)&\pnVar(Y)
 \end{pmatrix}\\
 &\quad=\pnE\left[
 \begin{pmatrix}X-\pnE[X]\\Y-\pnE[Y]\end{pmatrix}
 \begin{pmatrix}X-\pnE[X]&Y-\pnE[Y]\end{pmatrix}
 \right],
\end{align*}
其中矩阵的期望按元素计算。
\end{pnDefinition}

\begin{pnProposition}{协方差的可加性与和的方差}
有
\begin{align*}
 \pnCov(X,Y+Z)&=\pnCov(X,Y)+\pnCov(X,Z),\\
 \pnVar(X+Y)&=\pnVar(X)+\pnVar(Y)+2\pnCov(X,Y).
\end{align*}
\end{pnProposition}
\begin{proof}
利用期望的线性性，
\begin{align*}
 \pnCov(X,Y+Z)
 &=\pnE\bigl[(X-\pnE[X])\bigl(Y+Z-\pnE[Y]-\pnE[Z]\bigr)\bigr]\\
 &=\pnE\bigl[(X-\pnE[X])(Y-\pnE[Y])\bigr]\\
 &\quad+\pnE\bigl[(X-\pnE[X])(Z-\pnE[Z])\bigr].
\end{align*}
再展开 $\pnVar(X+Y)=\pnCov(X+Y,X+Y)$，并用
$\pnCov(X,Y)=\pnCov(Y,X)$，便得到第二个公式。
\end{proof}

\subsection{概率生成函数}

\begin{pnDefinition}{概率生成函数}
设随机变量 $X$ 取值于 $\mathbb N$，定义其概率生成函数为
\[
 G_X(z)=\sum_{x=0}^{\infty}p_X(x)z^x
 =p_X(0)+p_X(1)z+p_X(2)z^2+\cdots
 =\pnE[z^X].
\]
\end{pnDefinition}

可先把 $G_X$ 看作形式幂级数；作为复变量函数时，它在 $|z|\leq1$ 上绝对收敛。
在 $z=0$ 处按幂级数约定 $z^0=1$。

\begin{pnProposition}{归一化、导数与矩}
概率生成函数满足 $G_X(1)=1$。在 $0<|z|<1$ 内，逐项求导得到
\begin{align*}
 G_X'(z)&=\sum_{x\geq1}xp_X(x)z^{x-1}
          =\pnE[Xz^{X-1}],\\
 G_X''(z)&=\sum_{x\geq2}x(x-1)p_X(x)z^{x-2}
          =\pnE[X(X-1)z^{X-2}].
\end{align*}
其中系数为零的项按零处理。在 $z=1$ 处，将导数理解为从实数区间内趋近 $1$ 的左极限；
若所需矩有限，则
\begin{align*}
 G_X'(1)&=\pnE[X],\\
 G_X''(1)&=\pnE[X(X-1)]=\pnE[X^2-X],\\
 \pnE[X^2]&=G_X''(1)+G_X'(1),\\
 \pnVar(X)&=\pnE[X^2]-(\pnE[X])^2\\
 &=G_X''(1)+G_X'(1)-\bigl(G_X'(1)\bigr)^2.
\end{align*}
\end{pnProposition}

\begin{pnProposition}{独立随机变量之和的概率生成函数}
若取值于 $\mathbb N$ 的随机变量 $X,Y$ 独立，则
\[
 G_{X+Y}(z)=G_X(z)G_Y(z).
\]
\end{pnProposition}
\begin{proof}
对 $|z|\leq1$，由独立性有
\[
 G_{X+Y}(z)=\pnE[z^{X+Y}]
 =\pnE[z^Xz^Y]
 =\pnE[z^X]\pnE[z^Y]=G_X(z)G_Y(z).
\]
因此上述等式也给出相应形式幂级数的恒等式。
\end{proof}

\begin{pnRemark}{代入单位圆上的变量}
对实数 $t$，有
\[
 \pnE[\mathrm e^{\mathrm i tX}]=G_X(\mathrm e^{\mathrm i t}),
 \qquad \mathrm i^2=-1.
\]
\end{pnRemark}

\subsection{应用：等待连续两个正面}

\begin{pnExample}{连续出现两次正面时停止}
独立地重复抛一枚硬币，每次出现正面 $H$ 的概率为 $p$，
出现反面 $T$ 的概率为 $q=1-p$，其中 $0<p\leq1$。
当第一次连续出现两个正面时停止，令 $X$ 表示停止时已经抛掷的总次数。
因此 $X\in\{2,3,\ldots\}$，并且 $p>0$ 时试验以概率 $1$ 停止。

可能的停止序列包括
\[
 HH,\quad THH,\quad TTHH,\quad THTHH,\quad\ldots.
\]
每个停止序列都以 $HH$ 结尾，且在结尾之前没有出现过 $HH$。
\end{pnExample}

\paragraph{用带权序列表示概率生成函数}
记所有停止序列组成的集合为 $\mathcal S$。
对 $s\in\mathcal S$，记其长度为 $|s|$，正面、反面的个数分别为
$h(s),t(s)$。由各次抛掷的独立性，
\[
 \Pr[\text{停止序列为 }s]=p^{h(s)}q^{t(s)}.
\]
按停止时间或按停止序列求和，得到同一个生成函数：
\begin{align*}
 G_X(z)
 &=\sum_{x=2}^{\infty}\Pr[X=x]z^x\\
 &=\sum_{s\in\mathcal S}p^{h(s)}q^{t(s)}z^{|s|}.
\end{align*}
例如，序列 $THTHH$ 有 $3$ 个正面、$2$ 个反面，长度为 $5$，
所以它在生成函数中的贡献为 $p^3q^2z^5$。

若把每个字母的权重记为
\[
 W(H)=pz,\qquad W(T)=qz,
\]
则连接两个字串时权重相乘。例如
$W(HH)=p^2z^2$，$W(HT)=pqz^2$。
这样，$G_X(z)$ 就是所有停止序列的权重之和。

\paragraph{按开头分类：停止或重新开始}
根据前一至两次抛掷，可以把停止序列分成三类：
\begin{enumerate}
 \item 开头为 $HH$：第二次抛掷就停止，序列恰为 $HH$。
 \item 开头为 $T$：没有留下可接成 $HH$ 的末尾正面，从此重新开始等待。
 \item 开头为 $HT$：第二次出现反面，同样重新开始等待。
\end{enumerate}
后两类中，重新开始之后的抛掷与之前独立，规则也和最初相同。
因此有不交分解
\[
 \mathcal S=\{HH\}\mathbin{\dot\cup}T\mathcal S
             \mathbin{\dot\cup}HT\mathcal S,
\]
其中 $T\mathcal S=\{Ts:s\in\mathcal S\}$，
$HT\mathcal S=\{HTs:s\in\mathcal S\}$。
对各类的权重求和，得到
\[
 G_X(z)=p^2z^2+qzG_X(z)+pqz^2G_X(z).
\]
移项并解出 $G_X$：
\begin{equation}
 \boxed{G_X(z)=\frac{p^2z^2}{1-qz-pqz^2}
 =\frac{p^2z^2}{1-(1-p)z-p(1-p)z^2}.}
 \label{eq:pn-two-heads-pgf}
\end{equation}
由于分母的常数项为 $1$，上述关系也可以按形式幂级数理解。
代入 $z=1$ 时，$1-q-pq=p^2$，所以 $G_X(1)=1$，与归一化性质一致。

\paragraph{由生成函数求期望}
为避免直接对商求导，先把式~\eqref{eq:pn-two-heads-pgf} 写为
\[
 G_X(z)(1-qz-pqz^2)=p^2z^2.
\]
两边求导，得到
\[
 G_X'(z)(1-qz-pqz^2)+G_X(z)(-q-2pqz)=2p^2z.
\]
取 $z=1$，利用 $G_X(1)=1$ 及 $1-q-pq=p^2$，有
\begin{align*}
 p^2G_X'(1)-q-2pq&=2p^2,\\
 p^2G_X'(1)&=2p^2+q+2pq=1+p.
\end{align*}
因此
\[
 \boxed{\pnE[X]=G_X'(1)=\frac{1+p}{p^2}
 =\frac1{p^2}+\frac1p.}
\]
求方差时仍使用前面的通式
\[
 \pnVar(X)=G_X''(1)+G_X'(1)-\bigl(G_X'(1)\bigr)^2.
\]


\subsection{常见分布及其概率生成函数}

下面保留本次笔记中的分布和推导。所有概率质量函数在所列取值范围外均为 $0$。

\begin{pnExample}{Bernoulli 分布}
设 $0\leq p\leq1$。若 $X\sim\operatorname{Bern}(p)$，则
\[
 p_X(x)=\begin{cases}p,&x=1,\\1-p,&x=0,\end{cases}
 \qquad G_X(z)=1-p+pz.
\]
\end{pnExample}

\begin{pnExample}{二项分布}
设 $n\in\mathbb N_{>0}$，$0\leq p\leq1$。
若 $X\sim\operatorname{Bin}(n,p)$，则
\[
 p_X(x)=\binom nx p^x(1-p)^{n-x},\qquad x=0,1,\ldots,n.
\]
其概率生成函数为
\[
 G_X(z)=\sum_{x=0}^n\binom nx p^x(1-p)^{n-x}z^x
       =(1-p+pz)^n.
\]
这也对应于 $n$ 个相互独立的 $\operatorname{Bern}(p)$ 随机变量之和。
\end{pnExample}

\begin{pnExample}{有限集合上的均匀分布}
设 $S$ 为非空有限集合。若 $X\sim\operatorname{Unif}(S)$，则
\[
 p_X(x)=\begin{cases}1/|S|,&x\in S,\\0,&x\notin S.\end{cases}
\]
这里 $S$ 不必是整数集合，因此一般情形下不为它写概率生成函数。
\end{pnExample}

\begin{pnExample}{几何分布：两种计数约定}
考虑相互独立的抛硬币试验，每次出现正面的概率为 $p$，其中 $0<p\leq1$。

若 $T$ 表示第一次出现正面时的试验总次数，则
\[
 p_T(t)=(1-p)^{t-1}p,\qquad t\in\mathbb N_{>0},
\]
其概率生成函数为
\[
 G_T(z)=\sum_{t=1}^{\infty}p(1-p)^{t-1}z^t
       =\frac{pz}{1-(1-p)z}.
\]
若 $F$ 表示第一次出现正面之前连续出现反面的次数，则 $F=T-1$，且
\[
 p_F(f)=(1-p)^fp,\qquad f\in\mathbb N,
 \qquad G_F(z)=\frac{p}{1-(1-p)z}.
\]
两种约定的取值起点不同，使用几何分布时应说明是在计数总试验次数，还是此前的失败次数。
\end{pnExample}

\begin{pnExample}{负二项分布}
在上述独立试验中，设 $r\in\mathbb N_{>0}$，用 $N$ 表示第 $r$ 次出现正面之前
出现反面的总次数，记为 $N\sim\operatorname{NegBin}(r,p)$。
则
\[
 p_N(x)=\binom{r+x-1}{r-1}p^r(1-p)^x,\qquad x\in\mathbb N.
\]
最后一次试验必须为正面；此前 $r+x-1$ 次试验中恰有 $r-1$ 次正面，因而得到上述系数。
把每次正面之前的失败次数分别记为 $F_1,\ldots,F_r$，则这些随机变量相互独立，
均服从从 $0$ 开始计数的几何分布，且 $N=F_1+\cdots+F_r$。因此
\[
 G_N(z)=\left(\frac{p}{1-(1-p)z}\right)^r.
\]
\end{pnExample}

\begin{pnExample}{Poisson 分布与二项分布的极限}
设 $\lambda>0$。若 $X\sim\operatorname{Pois}(\lambda)$，则
\[
 p_X(x)=\frac{\mathrm e^{-\lambda}\lambda^x}{x!},\qquad x\in\mathbb N,
\]
其概率生成函数为
\[
 G_X(z)=\sum_{x=0}^{\infty}\frac{\mathrm e^{-\lambda}\lambda^x}{x!}z^x
       =\mathrm e^{\lambda(z-1)}.
\]
将固定规模的试验分为 $n$ 个小段，每段的发生概率取为 $\lambda/n$，相应地考虑
$X_n\sim\operatorname{Bin}(n,\lambda/n)$（取 $n\geq\lambda$）。其概率生成函数满足
\begin{align*}
 G_{X_n}(z)
 &=\left(1-\frac\lambda n+\frac\lambda n z\right)^n\\
 &=\left(1+\frac{\lambda(z-1)}n\right)^n
 \longrightarrow \mathrm e^{\lambda(z-1)}.
\end{align*}
这与 Poisson 分布的概率生成函数一致。
\end{pnExample}

\end{document}
